\documentclass[a4paper, noanswers, addpoints, 12pt]{exam}
\usepackage[T1]{fontenc}
\usepackage[utf8]{inputenc}
\usepackage{lmodern}
\usepackage[french]{babel}
\usepackage{xcolor}
\usepackage{listings}
\usepackage{graphicx}
\usepackage[most]{tcolorbox}
\usepackage{tikz}

% --- Configuration de la syntaxe Python ---
\definecolor{codegreen}{rgb}{0,0.6,0}
\definecolor{codegray}{rgb}{0.5,0.5,0.5}
\definecolor{codepurple}{rgb}{0.58,0,0.82}
\definecolor{backcolour}{rgb}{0.95,0.95,0.92}
 
\lstdefinestyle{mystyle}{
    backgroundcolor=\color{backcolour},   
    commentstyle=\color{codegreen},
    keywordstyle=\color{magenta},
    numberstyle=\tiny\color{codegray},
    stringstyle=\color{codepurple},
    basicstyle=\footnotesize\ttfamily,
    breakatwhitespace=false,         
    breaklines=true,                 
    captionpos=b,                    
    keepspaces=true,                 
    numbers=left,                    
    numbersep=5pt,                  
    showspaces=false,                
    showstringspaces=false,
    showtabs=false,                  
    tabsize=2
}
\lstset{style=mystyle, language=Python}

\title{Évaluation NSI - Architecture, OS et Système Linux}
\author{Babash - NSI Tani Malandi et Johanne Bakalara}
\date{Avril 2026, Durée: 1h30}

\usepackage[a4paper,top=2.3cm,bottom=2.8cm,left=2cm,right=2cm,headsep=0.6cm,footskip=1.2cm]{geometry}
\lstset{extendedchars=true,literate={â}{{\^a}}1}
\begin{document}
\runningfooter{}{\footnotesize Babash - NSI Tani Malandi et Johanne Bakalara \textperiodcentered{} CC BY-SA 4.0 \textperiodcentered{} page \thepage}{}
\firstpagefooter{}{\footnotesize Babash - NSI Tani Malandi et Johanne Bakalara \textperiodcentered{} CC BY-SA 4.0 \textperiodcentered{} page \thepage}{}

\maketitle
\thispagestyle{headandfoot}

\section*{Consignes:}
\begin{center} 
\fbox{\parbox{\linewidth}{
\begin{center}
La clarté de la rédaction et les explications font partie intégrante de la notation. 
Pour les questions demandant du code, respectez la syntaxe Python. 
\end{center}
}}
\end{center}

\section*{Barème :}
\begin{center}
  \par\bigskip
  \gradetable[h]
\end{center}

\vspace{0.5cm}

\begin{questions}

% --- SECTION 1 : ARCHITECTURE (2 points) ---
\section{Architecture Matérielle (2 points)}

\question[2] Complétez le tableau suivant en donnant la signification de l'acronyme et son rôle principal au sein de l'ordinateur :

\begin{center}
\renewcommand{\arraystretch}{2}
\begin{tabular}{|c|p{5cm}|p{7cm}|}
\hline
\textbf{Acronyme} & \textbf{Nom complet (Français ou Anglais)} & \textbf{Rôle principal} \\ \hline
CPU & & \\ \hline
GPU & & \\ \hline
RAM & & \\ \hline
ROM & & \\ \hline
\end{tabular}
\end{center}


% --- SECTION 2 : UNITÉS ET STOCKAGE (3 points) ---
\section{Unités et Stockage (3 points)}

\question[1] Pour le même prix et des performances équivalentes par ailleurs dans le commerce vaut-il mieux acheter une clé USB de 250Mbits ou une clé de 5Go. Vous justifierez par le calcul.

\question[2] On dispose d'un disque dur externe d'une capacité de \textbf{2 téra-octets (To)}.
\begin{parts}
    \part Combien de films pesant chacun \textbf{1 Go} peut-on stocker sur ce disque ? (Détaillez le calcul).
    \part On souhaite copier l'intégralité de ces films sur un autre serveur. Sachant que la vitesse d'écriture réelle du disque est de \textbf{20 Mo/s}, combien de temps faudra-t-il pour transférer \textbf{1 Go} de données ?
\end{parts}


% --- SECTION 3 : SYSTÈME D'EXPLOITATION (5 points) ---
\section{Système d'Exploitation (5 points)}

\question[2] Citez et expliquez brièvement les 4 rôles majeurs d'un Système d'Exploitation (OS) vus en cours.

\question[3] \textbf{Logiciel Libre vs Propriétaire :} 
\begin{parts}
    \part Quelle est la différence fondamentale concernant l'accès au "code source" entre un logiciel libre et un logiciel propriétaire ?
    \part Citez un exemple de système d'exploitation \textbf{libre} et un exemple de système \textbf{propriétaire}.
\end{parts}


% --- SECTION 4 : LINUX ET COMMANDES (10 points) ---
\section{Manipulation du Système Avalon (10 points)}

\begin{tcolorbox}[colback=white, colframe=black!75, title=Arborescence du serveur Avalon]
\centering
\begin{tikzpicture}[
    every node/.style={draw, rounded corners, fill=blue!5, font=\footnotesize, minimum width=2.2cm},
    edge from parent/.style={draw, -stealth, thick},
    level 1/.style={sibling distance=5cm, level distance=1.5cm},
    level 2/.style={sibling distance=2.5cm, level distance=1.5cm},
    level 3/.style={sibling distance=2.5cm, level distance=1.5cm}
]
\node {/ (racine)} 
    child {node {bin/ (D)}
        child {node {ls (F)}}
        child {node {cp (F)}}
    }
    child {node {home/ (D)}
        child {node {ghost/ (D)}
            child {node {note.txt (F)}}
            child {node {scripts/ (D)}
                child {node {test.sh (F)}}
            }
        }
        child {node {admin/ (D)}}
    }
    child {node {var/ (D)}
        child {node {log/ (D)}
            child {node {sys.log (F)}}
        }
    };
\end{tikzpicture}
\par\medskip
\small{\textbf{(D)} : Dossier \quad \textbf{(F)} : Fichier}
\end{tcolorbox}

\question[2] \textbf{Navigation :}
\begin{parts}
    \part Quelle commande permet d'afficher le chemin absolu du répertoire où l'on se trouve ? 
    \part Écrire la commande (chemin relatif) pour aller du dossier \texttt{bin} vers le dossier \texttt{admin}. \\
\end{parts}

\question[4] \textbf{Gestion de fichiers :} 
Écrire les commandes permettant de réaliser les actions suivantes depuis \texttt{/home/admin} :
\begin{parts}
    \part Créer un dossier nommé \texttt{credentials}.
    \part Déplacer le fichier \texttt{note.txt} dans le dossier \texttt{credentials}.
    \part Afficher le contenu du fichier \texttt{sys.log} (sans quitter votre répertoire).
    \part Supprimer le fichier \texttt{cp}.
\end{parts}

\question[4] \textbf{Droits d'accès (chmod) :} 
On observe le fichier \texttt{test.sh} avec la commande \texttt{ls -l} : \\
\texttt{-rw-r-{}-r-{}- 1 ghost users 1024 test.sh}

\begin{parts}
    \part Un utilisateur appartenant au groupe \texttt{users} peut-il exécuter ce script ? Justifiez.
    \part Donnez la valeur \textbf{octale} correspondant aux droits \texttt{rw-r-{}-r-{}-}.
    \part Donnez la commande permettant d'ajouter le droit d'écriture au groupe et aux "autres" (others)en mode symbolique.
\end{parts}

% --- SECTION BONUS : PYTHON (2 points) ---
\section{Bonus : Interprétation d'erreurs Python (2 points)}

\question[2] Pour chaque code ci-dessous, un message d'erreur est généré. Expliquez pourquoi cette erreur survient et proposez une correction du code.

\begin{lstlisting}[language=Python]
notes = [12, 15, 18]
for i in range(4):
    print(notes[i])
\end{lstlisting}
\begin{tcolorbox}[colback=red!5,colframe=red!75,title=Terminal Python, fonttitle=\bfseries]
\texttt{\color{red}IndexError: list index out of range}
\end{tcolorbox}

\begin{lstlisting}[language=Python]
age = input("Quel est votre âge ? ")
if age >= 18:
    print("Majeur")
else:
    print("Mineur")
\end{lstlisting}
\begin{tcolorbox}[colback=red!5,colframe=red!75,title=Terminal Python, fonttitle=\bfseries]
\texttt{\color{red}TypeError: '>=' not supported between instances of 'str' and 'int'}
\end{tcolorbox}

\begin{lstlisting}[language=Python]
def calcul(a, b):
    return a / b

print(calcul(10, 0))
\end{lstlisting}
\begin{tcolorbox}[colback=red!5,colframe=red!75,title=Terminal Python, fonttitle=\bfseries]
\texttt{\color{red}ZeroDivisionError: division by zero}
\end{tcolorbox}

\begin{lstlisting}[language=Python]
inventaire = {"pommes": 5, "bananes": 2}
print(inventaire["poires"])
\end{lstlisting}
\begin{tcolorbox}[colback=red!5,colframe=red!75,title=Terminal Python, fonttitle=\bfseries]
\texttt{\color{red}KeyError: 'poires'}
\end{tcolorbox}


\end{questions}

\end{document}